% \vspace{-10pt}
To address the challenges discussed in Section~\ref{sec-background}, we introduce \ourwork, a comprehensive framework that encapsulates a suite of machine learning components within user-friendly APIs. This framework is tailored for HPC users, simplifying the implementation process and making the robust capabilities of language models more accessible and user-friendly within the HPC community. The primary goal of \ourwork is to reduce the complexities inherent in employing language models, thus enabling HPC users to leverage their powerful capabilities more effectively and efficiently.




\subsection{\ourwork  design overview}
% DO NOT ADVERTIZE HUGGING FACE IN our own design overview
%While Hugging Face offers maturing designs for each step in a machine-learning application, a gap remains for HPC users looking to leverage these benefits. 
%To bridge this gap, we introduce \ourwork, a scalable and extensible framework specifically designed to effectively deploy language models for the HPC domain. 
%In line with the FAIR principles, \ourwork is built atop Hugging Face. This approach enables HPC users to easily access and deploy models and datasets shared by the NLP community. The APIs of Hugging Face and their documentation help HPC users reduce the steep learning curve typically associated with language models.

Figure~\ref{fig:architecture} provides the overview of the \ourwork framework. It is built on top of multiple internal machine learning components with \face-compatible APIs. Higher-level components provide concepts and interfaces to users, while middle or lower-level components provide implementation support. 
Table \ref{tab:api_table} shows the example classes and functions in LM4HPC API, including those supporting HPC-specific language models,  tokenizers for programming languages, datasets, inference pipelines, and evaluation.
%Figure~\ref{fig:architecture} provides an overview of the \ourwork framework. Delivered as a Python package, \ourwork aims to simplify the deployment of large language models within the HPC domain. The framework employs \face as the underlying runtime library to communicate between machine learning frameworks and higher-level APIs.
%The higher-level components within \ourwork comprise five primary parts: dataset, tokenizer, pipelines, vector stores, and evaluations. 
We elaborate on some essential components in the following subsections. 

\begin{figure}[h]
    \centering
    \includegraphics[width=\textwidth]{images/architecture.pdf}
    \caption{Overview of the \ourwork framework}
    \label{fig:architecture}
\end{figure}



% Please add the following required packages to your document preamble:
% \usepackage{multirow}
% \usepackage[table,xcdraw]{xcolor}
% If you use beamer only pass "xcolor=table" option, i.e. \documentclass[xcolor=table]{beamer}
\begin{table}[!t]
\centering
\caption{LM4HPC API: example classes and functions. Each class can be imported using ``from lm4hpc import *'' in Python.}
\label{tab:api_table}
\footnotesize 
\setlength{\tabcolsep}{2pt} 
\begin{tabular}{|p{0.13\columnwidth}|p{0.25\columnwidth}|p{0.56\columnwidth}|}
\hline
\textbf{\scriptsize LM4HPC classes} & \textbf{\scriptsize Description} & \textbf{\scriptsize Example API functions} \\
\hline
\multirow{3}{\linewidth}{\scriptsize \texttt{hpcmodel}} &
\multirow{3}{\linewidth}{\raggedright\arraybackslash \scriptsize Fine-tune text-based (HF, OpenAI) and graph-based models, including local private ones, for HPC tasks} &
\raggedright\arraybackslash \scriptsize \texttt{hpcmodel.from\_pretrained(model\_name\_or\_path: Optional[str], *model\_args, **kwargs)} \\ \cline{3-3} 
& & \raggedright\arraybackslash \scriptsize \texttt{hpcmodel.save\_pretrained(model\_name\_or\_path: str, *model\_args, **kwargs)} \\ \cline{3-3} 
& & \raggedright\arraybackslash \scriptsize \texttt{hpcmodel.finetune()} \\
\hline
\multirow{3}{\linewidth}{\scriptsize \texttt{hpctokenizer}} &
\multirow{3}{\linewidth}{\raggedright\arraybackslash \scriptsize APIs to represent code in either tokenized text, trees, or graphs} &
\raggedright\arraybackslash \scriptsize \texttt{hpctokenizer.from\_pretrained(model\_name\_or\_path: Optional[str], *model\_args, **kwargs)} \\ \cline{3-3} 
& & \raggedright\arraybackslash \scriptsize \texttt{hpctokenizer.addtokens(contentsingle\_word=False, strip=False, normalized=True)} \\ \cline{3-3} 
& & \raggedright\arraybackslash \scriptsize \texttt{hpctokenizer.encoding()} \\
\hline
\multirow{4}{\linewidth}{\scriptsize \texttt{hpcdatasets}} &
\multirow{4}{\linewidth}{\raggedright\arraybackslash \scriptsize Load and process HPC datasets} &
\raggedright\arraybackslash \scriptsize \texttt{hpcdatasets.load(path: str, data\_files: Union[str, List, Dict, None], **kwargs)} \\ \cline{3-3} 
& & \raggedright\arraybackslash \scriptsize \texttt{hpcdatasets.split(dataset: hpcdatasets, partition: [float, float, float], **kwargs)} \\ \cline{3-3} 
& & \raggedright\arraybackslash \scriptsize \texttt{hpcdatasets.shuffle(dataset: hpcdatasets, **kwargs)} \\ \cline{3-3} 
& & \raggedright\arraybackslash \scriptsize \texttt{hpcdatasets.sort(dataset: hpcdatasets, **kwargs)} \\
\hline
\scriptsize \texttt{hpcpipeline} &
\scriptsize Pre-built pipelines for common PLP tasks &
\raggedright\arraybackslash \scriptsize \texttt{hpcpipeline(task: str, model\_name\_or\_path: str, *model\_args, **kwargs)} \\
\hline
\multirow{2}{\linewidth}{\scriptsize \texttt{hpceval}} &
\multirow{2}{\linewidth}{\raggedright\arraybackslash \scriptsize Evaluate the results of various models} &
\raggedright\arraybackslash \scriptsize \texttt{hpceval.compute(task: str, models\_name\_or\_path: [[str]], data\_files: Union[str, List, Dict, None], *model\_args, **kwargs)} \\ \cline{3-3} 
& & \raggedright\arraybackslash \scriptsize \texttt{hpceval.plot(shape: str)} \\
\hline
\end{tabular}
\end{table}


% \vspace{-15pt}
%========================================
\subsection{HPC Tasks and Inference Pipelines}
HPC users are interested in a wide range of tasks related to programming language processing. Table~\ref{tab-HPC-tasks} outlines one way to categorize HPC-specific tasks. The purpose here is not to provide a comprehensive taxonomy of all tasks but a starting point for common tasks we are interested in supporting in our framework. Most tasks are self-explanatory by names and each may have further sub-tasks. For example, clone detection can be viewed as a specialized sub-task under code similarity analysis. 
\begin{table}[htbp]
\centering
\caption{HPC Tasks for Programming Language Processing: Categories and Examples}
\label{tab-HPC-tasks}
\begin{tabular}{l|l|l}
\hline
%\multicolumn{3}{c}{\textbf{Category}} \\ \hline
\textbf{Code Analysis} & \textbf{Code Generation} & \textbf{Others} \\ \hline
Compiler Analysis & Code Completion & Test Case Generation \\
Algorithm Classification & Natural Language-to-Code & Code Search \\
Code Similarity Analysis & Code Translation & Question Answering \\
Documentation Generation & Code Repair & Code Review \\
Parallelism Detection & Code Migration & Decompilation  \\
Defect Detection & Compilation &  IR-to-Source Translation \\ \hline
\end{tabular}
\end{table}

In the context of machine learning, a pipeline represents a sequence of data processing stages to complete a task. Our \ourwork framework extends the pipeline function provided by \face, adapting it for HPC tasks.  
%Pipelines allow HPC researchers and developers, even those with limited ML experience, to rapidly harness the power of cutting-edge LLMs. Currently, \face provides pipelines for NLP and image processing tasks only. 
%\todo{HPCFAIR vs. HPC-FAIR: which one do you mean?  This paper focuses on HPC-FAIR. We only need to briefly mention previous HPCFAIR in related work. }
%In an effort to streamline this process and provide user-friendly interfaces for executing a variety of HPC-related inference tasks using pre-trained models, we have extended the \face pipeline() API within \ourwork. 
We have developed three inference pipelines: code similarity analysis, parallelism detection, and OpenMP question answering. Code similarity analysis determines the similarity between a pair of code snippets. 
%The exact meaning of similarity depends on the context. It could mean similar syntax, semantics, or functionality. 
Parallel detection is defined to check if an input code snippet can be parallelized or not using OpenMP.  The OpenMP question answering pipeline is designed to use models to generate answers to OpenMP-related questions. %\todo{Double check if the task description is accurate} %These pipelines are designed to significantly reduce the complexity of integrating machine learning models into HPC workflows.
% \begin{comment}  % No space   

% Listing \ref{ls:pipepline} illustrates how to carry out inference tasks using \ourwork. In line 3, the parallelism detection pipeline is loaded by passing the pipeline identifier to the \mintinline{python}{pipeline()} API. 
% With just two lines of code, the pipeline API can generate results for parallelism detection tasks. Each pipeline task is equipped with predefined pre-trained models and tokenizers. This efficient design drastically reduces the complexity of these tasks, enabling users to obtain and evaluate results quickly and effortlessly.

% \begin{listing}[h]
% \caption{Example of using \ourwork pipeline.}
% \label{ls:pipepline}
% \begin{minted}
% [
% frame=lines,
% framesep=2mm,
% baselinestretch=1.2,
% % bgcolor=LightGray,
% fontsize=\footnotesize,
% linenos
% ]{python}
% from \ourwork import pipeline
% # parallelism detection task
% pipe = pipeline("parallelism-detection")
% res = pipeline(code)
% \end{minted}
% \end{listing}
% \end{comment}
% we have to mention pipelines before Datasets. Otherwise hard to explain datasets internals like enabling parallel detection, similarity analysis etc.

Tokenizers are responsible for preprocessing input into an array of numbers as inputs to a model. They are essential components used by pipelines. Most LM tokenizers are primarily designed for NLP tasks. 
%Consequently, they often fall short when directly applied to HPC tasks. 
For instance, given a function name \texttt{my\_func}, a typical NLP tokenizer like BERT might split it into separate tokens (such as `\texttt{my}', `\texttt{\_}', and `\texttt{func}') while a code-aware tokenizer may treat the function name as a single entity to ensure a more meaningful representation. % of the code. %This clearly demonstrates the limitations of NLP tokenizers in correctly parsing code.

To overcome this, we developed the \ourwork tokenizer, leveraging the treesitter~\cite{treesitter} and programl~\cite{programl} library. Our tokenizer is specifically designed to handle the pre-processing of code data required for a language model. It includes tokenizers such as the ast-tokenizer.
%which parses code into tokens based on Abstract Syntax Tree (AST) traversal, thereby ensuring a more accurate representation of the code. 
As a result, \ourwork can accommodate models (such as augAST~\cite{chen2023learning-dataset}) that require AST as input in the pipeline. % and DeepSCC~\cite{}.
%augAST 

\subsection{Datasets}
\label{sec-datasets}
Datasets are crucial for any machine learning application. 
%The NLP community shares datasets via \face, establishing a platform to evaluate the performance of different models using the same dataset. 
Within the \ourwork framework, we contribute HPC-specific datasets either by converting existing ones into \face-compatible formats or by creating new ones from scratch. 

%in two significant ways: %1. We convert existing HPC datasets into a \face compatible format and share them on the \face Hub. This is achieved by transforming the HPC data into a JSON file where each row represents a single example from the HPC dataset. 
% 2. We create new HPC datasets specifically designed for LLMs. In this work, we have processed three datasets and evaluated different models with these datasets in section 4. 
    %1) We convert existing HPC datasets into a \face-compatible format and share them on the \face Hub. This is achieved by transforming the HPC data into a JSON file where each row represents a single example from the HPC dataset. 
    %2) We create new HPC datasets.     
%In this work, we have processed three datasets and evaluated different models with these datasets in section 4. 


We have converted three existing datasets to be compatible with \face dataset API: POJ-104~\cite{dataset:poj104}, DRB-ML~\cite{osti_1958879}, and \crawldata~\cite{chen2023learning-dataset}.  
POJ-104 is derived from a pedagogical programming open judge (OJ) system that automatically evaluates the validity of submitted source code for specific problems by executing the code. This dataset is particularly useful for the code similarity task. % as it consists of 104 problems, each accompanied by 500 student-written C/C++ programs that provide solutions to the respective problem. 
The DRB-ML dataset contains 658 C/C++ OpenMP kernels derived from DataRaceBench~\cite{dataset:dataracebench}. We extended it to have labels indicating if a kernel is parallelizable or not. 
The \crawldata dataset is an open-source benchmark composed of data from three resources: code crawled from GitHub, OpenMP benchmarks such as Nas Parallel Benchmarks~\cite{jin1999openmp} and Rodinia~\cite{che2009rodinia}, and synthetic code. This dataset contains loops with labels indicating whether a loop is parallel and, if parallelizable, the corresponding OpenMP directive associated with the loop. 

Furthermore, we have manually created a new OpenMP question answering dataset called OMPQA in order to probe the capabilities of language models in single-turn interactions with users. Similar to other QA datasets, we include some request-response pairs which are not strictly question-answering pairs. 
%We manually 
%aim to create a comprehensive, diverse, and balanced OpenMP Q\&A dataset. 
The categories and examples of questions in the OMPQA dataset can be found in  Table~\ref{tab-omp-qa}. 
%The basics category include question-answer pairs related to basic OpenMP concepts. The example category collects questions about code examples written in OpenMP. The compilers category list questions regarding OpenMP compilers, including their usage and internals. The benchmarks category has questions probing the knowledge about popular OpenMP benchmarks for performance, correctness, and so on.


\begin{table}[h]
\centering
\begin{tabular}{|c|c|p{3.7in}|}
\hline
 \textbf{Category}&  \textbf{Count} & \textbf{Example Questions}    \\
\hline
 Basics & 40 &  What is a worksharing construct in OpenMP? \\
\hline
Examples & 20 & Give an example OpenMP C code for computing PI using numerical integration. \\
\hline
Compilers & 24 &  {In what language is LLVM written? \newline How is a parallel region represented in Clang?} \\
\hline
 Benchmarks & 23 & {What are the NAS Parallel benchmarks? \newline Which benchmark assesses data race detection tools?}  \\
\hline 
\end{tabular}
\caption{OMPQA: categories and examples of questions}
% current collection of Q&A
% https://docs.google.com/spreadsheets/d/1dqsd2v_lcRyLv8XiLIy00E_ezBiaivcSjZrMDqFeu_w/edit#gid=0
\label{tab-omp-qa}
\end{table}


\begin{comment}
    
%\subsection{Tokenizer} % merged into the discussion of pipelines
Most LM tokenizers are primarily designed for NLP tasks. Consequently, they often fall short when directly applied to HPC tasks. For instance, given a function name \mintinline[]{python}{my_func}, an NLP tokenizer such as BERT might split it into separate tokens [\mintinline[breaklines, breakafter=]{bash}{'my', '_', 'func'}]. This clearly demonstrates the limitations of NLP tokenizers in correctly parsing code.

To address this issue, we have developed the \ourwork tokenizer, built on top of various code tokenizer libraries. The \ourwork tokenizer is specifically designed to handle all the pre-processing of code data required for a language model. It includes tokenizers such as the ast-tokenizer, which parses code into tokens based on Abstract Syntax Tree (AST) traversal, thereby ensuring a more accurate representation of the code.

Listing \ref{ls:hpcfair-tokenizer} \todo{Where is the listing? We don't have much space, we may just briefly mention tokenizers  as part of the pipeline}demonstrates the steps involved in tokenizing code and processing it through models. This procedure can be complex for those with limited machine learning backgrounds. Challenges such as selecting the appropriate model, choosing the right tokenizer, and managing post-processing may demand substantial effort to address. 
\end{comment}

\vspace{-25pt}
\subsection{Integration With New Data}
% https://drive.google.com/file/d/17iV8BbkY5QCwT3G6r6zfwQaOpYanWjV4/view?usp=sharing
% LangChain source code

% https://docs.google.com/drawings/d/10c3a9qfFxfLctr9KiROoK2xy5NhLFfW2TMg3cr_S-x4/edit?usp=sharing
% langchain figure if space allows
%Large language models are trained on a variety of data and have demonstrated their competence in a wide range of NLP tasks. Remarkably, they can also handle some HPC tasks when their training data contains programming languages. 
%However, these out-of-the-box models may not always provide the optimal performance for more complex HPC tasks. 
Language models derive knowledge from training datasets and store this knowledge in internal weights within the model's neural network architecture. However, incorporating new information into a trained model presents a challenge. Traditionally, one might fine-tune pre-trained models with their own data for specific tasks, but this approach requires substantial relevant data and can be resource-intensive. An alternative approach involves integrating new data as context information into a user prompt, but this is constrained by the limited input token lengths of current models.

To address this challenge, \ourwork leverages the LangChain framework~\cite{langchain2023} to easily integrate new data. LangChain aggregates a wide variety of components to build applications using LLMs. Particularly, it provides APIs allowing LLM applications to store large amounts of text in semantic databases called vector stores. 
The way to integrate new data can be done in two steps. First, text data is chunked and embedded with an LLM before being saved into a vector store. 
%Later, The embeddings and the corresponding text chunks are then retrieved using similarity in the latent space. 
Later, user prompts are matched with relevant chunks in the vector store using similarity analysis. 
%The matching chunks are injected in the prompt to provide context.
The top-matched chunks are then injected into the original prompts to form a new prompt with relevant context information. By employing this new prompt, language models can generate answers that incorporate new and relevant user data while still staying within the token length limits.
%generate Vector Store from their codebase or documentation. Initial queries from users are sent to the Vector Store for a similarity search instead of being sent directly to the LLMs. Consequently, LLMs receive both the initial query and the relevant information from the users' database, enabling them to generate more accurate answers or actions.
%Our solution with Langchain vector storage is designed to integrate user data with LLMs, offering a more convenient and efficient approach to tailor the LLMs to the specifics of users' HPC tasks, thereby enhancing their performance.

\vspace{-10pt}
\subsection{Evaluation}
An easily accessible harness for evaluating different language models on HPC tasks is crucial. Standard and reproducible results from such evaluations can provide researchers and developers with insightful starting points, helping them select suitable models for their specific needs and identify research opportunities. 

In response to this need, we developed an evaluator API in \ourwork. One challenge we encountered is the lack of standardized metrics for code evaluation. Unlike natural language tasks, where metrics such as BLEU,
% ~\cite{papineni2002bleu}, 
ROUGE,
% ~\cite{lin2004rouge}
and METEOR
% ~\cite{banerjee2005meteor} 
are commonly used, the domain of code lacks such universally accepted measures of quality. We are adding various LLM metrics such as CodeBLEU~\cite{ren2020codebleu} for code output.
%These metrics aim to provide a more comprehensive and accurate evaluation of models' performance on code-related tasks. 
Another challenge is that language models may generate different answers for the same input in different inference runs. Evaluation should consider consistent sampling settings (such as temperatures) and control over random seeds to improve reproducibility.  


%By integrating these metrics into \ourwork, we hope to facilitate more robust, effective, and straightforward model evaluation for HPC tasks, ultimately aiding in the selection and development of more advanced models for these complex tasks.
%\todo{\textcolor{blue}{CodeBLEU} what new metrics?}

%\todo{\textcolor{blue}{added }mention that we create and release a set of test harnesses (scripts) to enable standard and reproducible evaluation of supported HPC tasks.}

Ultimately, many users are interested in seeing leaderboards that showcase mainstream models competing on common HPC tasks. To satisfy this interest, we create and release a set of test harnesses scripts to enable standard and reproducible evaluation for supported HPC tasks.
% With its comprehensive components and APIs, our framework supports the generation of such leaderboards using a set of test harness scripts. 
%This allows the community to easily identify the most appropriate models for a given task, as well as the areas that require further research and development.  

%\todo{\textcolor{blue}{table 2} We miss the table of the APIs we provide and what functionality each of them provides}


\vspace{-10pt}